% =============================================================================
\chapter{Testes e reprodutibilidade}\label{cap:reprodutibilidade}
% =============================================================================

\section{Testes automatizados}\label{sec:testes}

O arquivo \code{tests/test_etl_and_modeling.py} reúne sete testes escritos com
pytest. Seis deles usam pequenos conjuntos sintéticos e não dependem da base
real; o sétimo executa o painel e só roda quando os Parquets processados estão
disponíveis (\autoref{qua:testes}).

\begin{quadro}[htbp]
\caption{\label{qua:testes}Testes automatizados do projeto}
\begin{corpotabela}
\begin{tabularx}{\textwidth}{|L{6.4cm}|Y|}
\hline
\textbf{Teste} & \textbf{O que garante} \\ \hline
\code{test_clean_and_validate_trades_contract} & A limpeza remove \code{ingest_ts} inválido e as colunas constantes, preserva \code{exchange} e \code{symbol} nos metadados, mapeia corretamente o lado do agressor e produz uma base que passa no contrato de qualidade \\ \hline
\code{test_build_bars_and_fill_bar_gaps_no_lookahead} & A grade de 1 segundo tem o número correto de barras; segundos vazios aparecem como nulos e são preenchidos com o preço \textbf{anterior}, nunca com o futuro; volume zero, OFI neutro e VWAP igual ao último preço \\ \hline
\code{test_build_taker_orders_merges_consecutive_same_ts_and_side} & Negócios consecutivos de mesmo milissegundo e lado formam uma única ordem, com quantidade somada e \textit{book walk} positivo \\ \hline
\code{test_build_lagged_features_strictly_causal_no_lookahead} & Perturbar drasticamente preços, OFI e volume a partir de um instante não altera nenhum atributo anterior a ele \\ \hline
\code{test_daily_walk_forward_folds_chronological} & Todas as observações de treino são anteriores ao dia de teste, e cada dobra testa exatamente um dia \\ \hline
\code{test_prepare_modeling_dataset_shapes} & O conjunto de modelagem não tem nulos, atributos e alvos têm o mesmo tamanho, e os grupos de atributos existem nas colunas \\ \hline
\code{test_streamlit_dashboard_renders_without_errors} & O painel renderiza sem exceções, com cinco indicadores e quatro abas (ignorado se os Parquets não existirem) \\ \hline
\end{tabularx}
\end{corpotabela}
\fonteautores[Elaborado]
\end{quadro}

Para esta documentação, a suíte foi executada em 25/11/2026, em ambiente limpo,
com Python 3.13, pandas 3.0, NumPy 2.5, scikit-learn 1.9 e Streamlit 1.67. O
resultado foi de \textbf{seis testes aprovados e um ignorado}: o teste do painel
foi ignorado porque, nesse ambiente, os Parquets processados não haviam sido
gerados.

\section{Como executar}\label{sec:comoexecutar}

\subsection{No Google Colab}\label{sec:colab}

É a forma recomendada para o notebook do \textit{Checkpoint}~2, que baixa os
dados sozinho: no Colab, escolher \textit{Arquivo} $\rightarrow$ \textit{Fazer
upload de notebook}, selecionar o arquivo \code{.ipynb} e então
\textit{Ambiente de execução} $\rightarrow$ \textit{Executar tudo}. A execução
leva de 1 a 3 minutos e cabe na memória do Colab gratuito.

\subsection{Localmente}\label{sec:local}

A partir da raiz do repositório, criar um ambiente virtual e instalar as
dependências:

\begin{lstlisting}[language=bash]
python -m venv .venv && source .venv/bin/activate
pip install -r requirements.txt
\end{lstlisting}

Em seguida, gerar e validar os Parquets de \code{data/processed/} pelo pipeline
consolidado (a função é idempotente: se os arquivos já existirem, nada é
refeito):

\begin{lstlisting}[language=bash]
python -m src.etl
\end{lstlisting}

Os notebooks são abertos a partir da pasta \code{notebooks/}, o painel é iniciado
como descrito na \autoref{sec:dashexecucao}, e os testes rodam a partir da raiz:

\begin{lstlisting}[language=bash]
cd notebooks && jupyter notebook A3_CP3_modelagem_storytelling.ipynb
streamlit run app/dashboard.py
python -m pytest -q tests/
\end{lstlisting}

\section{Requisitos de ambiente e cuidados com memória}\label{sec:memoria}

A base tem 27,9 milhões de linhas, e o Colab gratuito oferece cerca de 12~GB de
RAM. O pico de uso medido foi de aproximadamente 6~GB com pandas~3 e 9~GB com
pandas~2.2. Para caber nesse limite, o projeto adota as seguintes técnicas:

\begin{alineas}
  \item textos armazenados no formato compacto do Arrow (padrão no pandas 3 e
        ativado explicitamente no 2.x);
  \item alocador de memória do sistema no PyArrow, que devolve ao sistema
        operacional a memória liberada, e liberação explícita após leituras
        grandes;
  \item transformação sobre o próprio \textit{DataFrame} carregado, com o bruto
        preservado em disco, e remoção de objetos auxiliares seguida de coleta de
        lixo;
  \item leitura apenas das colunas necessárias nos produtos derivados e
        \textit{cache} das agregações usadas pelo painel.
\end{alineas}

A correlação de Spearman no pandas depende da biblioteca SciPy, que não aparece
explicitamente em \code{requirements.txt}, mas é instalada como dependência do
scikit-learn e já vem no Colab.

\section{Garantias de reprodutibilidade}\label{sec:garantias}

Os resultados podem ser reproduzidos por qualquer pessoa porque o projeto fixa
todas as fontes de variação controláveis:

\begin{alineas}
  \item \textbf{dados}: \textit{commit} da base e SHA-256 fixados; candles da
        Binance conferidos pelo SHA-256 publicado pela própria corretora;
  \item \textbf{código}: um único pacote (\code{src/}) compartilhado por
        notebooks, painel e testes, versionado em Git;
  \item \textbf{aleatoriedade}: semente fixa (\code{random_state=42}) nos modelos
        de árvore; os demais modelos são determinísticos;
  \item \textbf{ordem temporal}: dobras \textit{walk-forward} cronológicas e
        atributos estritamente causais, verificados por testes.
\end{alineas}
